\chapter{Instrumental Variables}\label{ch:b3:instrumental-variables}

Countries that trade more are richer. In 2019, across 157 countries, a
least-squares line through the logarithm of output per worker and the
logarithm of the trade share says that ten per cent more trade comes with
about eight per cent more output. % output: ex08_no_population (OLS 0.833)
Nobody can read that slope as the effect of trade. Rich countries can afford
ports, roads and the institutions that make trading cheap, so income may
cause trade as much as trade causes income, and whatever makes a country
productive may make it open too. The Year 2 volume met the same trap in a
single equation and called it omitted-variable bias. This chapter is about
the oldest way out of it: find something that moves the cause and touches the
outcome through no other channel. For trade, that something can be geography.
A country cannot choose its distance from its neighbours, and distance,
population and shared borders predict a good part of who trades with whom.

\begin{openingchart}{A first stage: the trade share a gravity regression
predicts from geography and population alone (horizontal), against the
actual trade share in goods and services (vertical), 2019, 157 countries,
both in per cent of GDP, on logarithmic scales. The line is the least-squares
fit of one logarithm on the other.}{CEPII Gravity database (BACI flows,
population-weighted distance, contiguity, population); World Bank, World
Development Indicators (NE.TRD.GNFS.ZS); Penn World Table 11.0 (country
coverage); our gravity fit}
\begin{tikzpicture}
\begin{loglogaxis}[omscatter, xlabel={Predicted trade share from geography (\% of GDP, log scale)},
  ylabel={Actual trade share (\% of GDP, log scale)}, xmin=1, xmax=1000, ymin=10, ymax=1000,
  xtick={1,10,100,1000}, ytick={10,100,1000}, log ticks with fixed point,
  height=0.6\linewidth]
  \addplot table[x=built, y=trade] {figdata/bachelor-3/01-instrumental-variables/out/first_stage.csv};
  \addplot[omThm, thick, no marks, sharp plot] table[x=built, y=trade] {figdata/bachelor-3/01-instrumental-variables/out/first_stage_line.csv};
  \node[font=\footnotesize, anchor=east] at (axis cs:470,80) {Congo};
  \node[font=\footnotesize, anchor=west] at (axis cs:2.4,118) {Fiji};
  \node[font=\footnotesize, anchor=west] at (axis cs:21,330) {Singapore};
  \node[font=\footnotesize, anchor=north] at (axis cs:5.1,25) {United States};
\end{loglogaxis}
\end{tikzpicture}
\end{openingchart}

Each dot is a country. The horizontal position is not a measurement of trade:
it is what a regression that knows only distances, borders and populations
predicts that country's trade to be. The vertical position is the trade the
country actually did. The cloud slopes up. A country whose geography predicts
ten per cent more trade trades, on average, about 2.4 per cent more
% figdata: first_stage_stats (slope 0.235)
(slope 0.235, standard error 0.058), and the $F$ statistic of that slope is
16.7. % figdata: first_stage_stats
The cloud is wide: Fiji, far from everyone, trades more than its geography
predicts, and Singapore and Hong Kong trade far more, because they re-export
what passes through their ports. Congo is the other extreme: its capital faces
Kinshasa across a river, so its predicted share is huge.
% figdata: first_stage (COG 490, FJI 2.3, SGP 20, USA 5.1)
That upward slope is the first of two things an instrument needs. The second
cannot be seen on any chart, and most of this chapter is about it.

\begin{recall}
The Year 2 volume wrote the multiple regression $y = X\beta + u$, solved the
normal equations, $\hat\beta = (X'X)^{-1}X'y$, and proved that OLS is
consistent when $\E[u \mid X] = 0$. When a regressor that matters is left in
the error, $\plim \hat\beta_1 = \beta_1 + \gamma\,\Cov(x, a)/\Var(x)$:
omitted-variable bias. It also defined potential outcomes $Y_i(1)$ and
$Y_i(0)$, the average treatment effect, selection bias, natural experiments,
classical measurement error and its attenuation bias, the $F$ statistic and
heteroskedasticity-robust standard errors. The Year 1 volume posed the
identification problem: a cloud of prices and quantities is not a curve.
\end{recall}

\begin{notation}[Symbols of this chapter]\label{not:b3:instrumental-variables:symbols}
In this chapter $z$ is an instrument and $Z$ the matrix of instruments, $\ell$
the number of columns of $Z$, and $W$ the matrix of exogenous controls (not a
nominal wage). A first-stage coefficient is written $\pi$, and $\pi_c$, $\pi_d$
are shares of compliance types: none of them is inflation. $D_i$ is a binary
treatment (in a diagram, the curve $D$ is still a demand curve); $s$ is years
of schooling and $a$ ability. In the supply-and-demand example, $\alpha_0$,
$\alpha_1$, $\gamma_0$ and $\gamma_1$ are the curves' coefficients, not a capital share.
\end{notation}

\section{Endogenous regressors}\label{sec:b3:instrumental-variables:endogenous}

Throughout this section the model is the simple linear regression
\begin{equation}\label{eq:b3:iv:model}
  y_i = \beta_0 + \beta_1 x_i + u_i, \qquad i = 1, \dots, n,
\end{equation}
with independent and identically distributed draws. The coefficient $\beta_1$
is a causal effect: the change in $y$ that a unit change in $x$ would cause
with $u$ held fixed. The question is whether OLS recovers it.

\begin{definition}[Endogenous regressor]\label{def:b3:instrumental-variables:endogenous}
\index{endogenous regressor}
In model~\eqref{eq:b3:iv:model}, the regressor $x$ is an \emph{endogenous
regressor} when $\Cov(x_i, u_i) \neq 0$, and an exogenous regressor when
$\Cov(x_i, u_i) = 0$.
\end{definition}

The word is borrowed from the Year 1 volume's models, where an endogenous
variable is one the model solves for. Here it means something narrower: a
regressor correlated with the error. The two senses meet in the first of three
classic cases.

\begin{enumerate}
\item \textbf{Simultaneity.} In a market, the price and the quantity are both
solved for by supply and demand. A regression of quantity on price mixes the
two curves, which is the Year 1 volume's identification problem in a
regression's clothes. The end of \cref{sec:b3:instrumental-variables:tsls}
draws it.
\item \textbf{Omitted variables.} A worker's ability raises both her schooling
and her wage. If ability is not measured, it sits in $u$, and schooling is
correlated with $u$. This is ability bias.
\item \textbf{Measurement error.} If schooling is recorded with a classical
error, the recorded value is correlated with the error of the regression on
the recorded value, and OLS is attenuated towards zero.
\end{enumerate}

In each case OLS converges to the wrong number. With $\hat\beta_1 =
\widehat{\Cov}(x, y)/\widehat{\Var}(x)$,
\[
  \plim \hat\beta_1^{\mathrm{OLS}} = \beta_1 + \frac{\Cov(x, u)}{\Var(x)},
\]
and the second term is the bias. No amount of data removes it: a larger sample
only estimates the wrong number more precisely.

\subsection*{A simulated town}

A simulation shows the size of the problem when the truth is known. In
\cref{lst:b3:instrumental-variables:data}, ability $a$ is a standard normal
variable that the economist never sees. Schooling is twelve years, plus one
year per standard deviation of ability, plus $0.8$ of a year for people who
grew up near a college, plus noise. The log wage rises by $0.08$ per year of
school, the true return, and by $0.10$ per standard deviation of ability.

\omcode{code/bachelor-3/01-instrumental-variables/listings/schooling_data.py}{1}{12}{The simulated town. Ability is drawn, used, and then hidden from the regressions.}\label{lst:b3:instrumental-variables:data}

\omcode{code/bachelor-3/01-instrumental-variables/listings/ability_bias.py}{1}{16}{OLS on the simulated town, and the omitted-variable formula.}\label{lst:b3:instrumental-variables:ability}

\omoutput{code/bachelor-3/01-instrumental-variables/outputs/ability_bias.txt}{OLS overstates the return by more than half.}\label{out:b3:instrumental-variables:ability}

The OLS slope is 0.1259 against a true 0.08. % output: ability_bias
The formula of the Year 2 volume predicts the gap. The variance of schooling
is $0.8^2 \times 0.25 + 1 + 1 = 2.16$ and its covariance with ability is 1, so
the limit of OLS is $0.08 + 0.10/2.16 \approx 0.126$. A regression with more
controls would not help unless one of them were ability itself. What helps is
a variable that moves schooling for reasons that have nothing to do with
ability. Growing up near a college is such a variable in this town, by
construction.

\section{The instrument and the Wald estimator}\label{sec:b3:instrumental-variables:instrument}

\begin{definition}[Instrumental variable]\label{def:b3:instrumental-variables:instrument}
\index{instrumental variable}\index{relevance condition}\index{exclusion restriction}
In model~\eqref{eq:b3:iv:model}, a variable $z$ is an \emph{instrumental
variable} (an instrument) for $x$ when it satisfies two conditions:
\begin{enumerate}
\item the \emph{relevance condition}, $\Cov(z_i, x_i) \neq 0$;
\item the \emph{exclusion restriction}, $\Cov(z_i, u_i) = 0$: the instrument
is uncorrelated with every other cause of $y$, so it affects $y$ only through
$x$.
\end{enumerate}
\end{definition}

Relevance is a statement about two observed variables and can be checked; the
opening chart checks it for geography. The \omterm{def:b3:instrumental-variables:instrument}{exclusion restriction} is a
statement about $u$, which is never observed. It is argued, not tested, and
every application of the method stands or falls on that argument. The
argument has two halves. The instrument must be as good as randomly assigned,
so that it is not correlated with the unobserved causes; and it must have no
path to $y$ of its own.

\begin{omfigure}
\begin{tikzpicture}[econ, every node/.style={font=\small},
  var/.style={draw, thick, circle, minimum size=10mm}]
  \node[var] (z) at (0,0) {$z$};
  \node[var] (x) at (3,0) {$x$};
  \node[var] (y) at (6,0) {$y$};
  \node[var, dashed] (u) at (4.5,2) {$u$};
  \draw[->, thick] (z) -- node[below] {relevance} (x);
  \draw[->, thick] (x) -- node[below] {$\beta_1$} (y);
  \draw[->, thick, dashed] (u) -- (x);
  \draw[->, thick, dashed] (u) -- (y);
  \draw[->, thick, omDebate, dotted] (z) to[bend right=40] node[below] {no direct path} (y);
  \draw[omDebate, thick] (2.85,-1.25) -- (3.15,-0.95) (2.85,-0.95) -- (3.15,-1.25);
  \draw[omDebate, thick, dotted] (z) -- node[above left] {no link} (u);
  \draw[omDebate, thick] (2.1,0.85) -- (2.4,1.15) (2.1,1.15) -- (2.4,0.85);
\end{tikzpicture}
\par\smallskip{\small The instrument $z$ moves $x$; the unobserved $u$ moves
both $x$ and $y$. The \omterm{def:b3:instrumental-variables:instrument}{exclusion restriction} removes the two crossed links: $z$
is unrelated to $u$, and $z$ reaches $y$ only through $x$.}
\end{omfigure}

\begin{proposition}[Consistency of the instrumental-variables estimator]\label{prop:b3:instrumental-variables:consistency}
Let $z$ be an instrument for $x$ in model~\eqref{eq:b3:iv:model}, and let
\[
  \hat\beta_1^{\mathrm{IV}} = \frac{\sum_i (z_i - \bar z)(y_i - \bar y)}{\sum_i (z_i - \bar z)(x_i - \bar x)}.
\]
Then $\plim \hat\beta_1^{\mathrm{IV}} = \beta_1$. If the \omterm{def:b3:instrumental-variables:instrument}{exclusion restriction}
fails, $\plim \hat\beta_1^{\mathrm{IV}} = \beta_1 + \Cov(z, u)/\Cov(z, x)$.
\end{proposition}

\begin{proof}
Substitute $y_i - \bar y = \beta_1 (x_i - \bar x) + (u_i - \bar u)$ in the
numerator:
\[
  \hat\beta_1^{\mathrm{IV}} = \beta_1 + \frac{\frac1n \sum_i (z_i - \bar z)(u_i - \bar u)}{\frac1n \sum_i (z_i - \bar z)(x_i - \bar x)}.
\]
By the law of large numbers the sample covariances converge in probability to
$\Cov(z, u)$ and $\Cov(z, x)$. The second is not zero by relevance, so the
ratio converges to $\Cov(z, u)/\Cov(z, x)$ (a continuous function of the two
limits). It is zero under the \omterm{def:b3:instrumental-variables:instrument}{exclusion restriction}.
\end{proof}

The second sentence of the proposition is the one to remember. A failure of
the \omterm{def:b3:instrumental-variables:instrument}{exclusion restriction}, $\Cov(z, u)$, is divided by the strength of the
\omterm{def:b3:instrumental-variables:tsls}{first stage}, $\Cov(z, x)$. With a strong instrument a small violation does
small harm. With a weak one, the same small violation is magnified:
\cref{sec:b3:instrumental-variables:weak} returns to it.

When the instrument takes two values, as an offer, a lottery or a birth date
before or after a cut-off does, the estimator has a simple form.

\begin{definition}[Wald estimator]\label{def:b3:instrumental-variables:wald}
\index{Wald estimator}
When the instrument $z$ is binary, the \emph{Wald estimator} is the ratio of
the difference in mean outcomes to the difference in mean regressors between
the two groups:
\[
  \hat\beta^{\mathrm{W}} = \frac{\bar y_1 - \bar y_0}{\bar x_1 - \bar x_0},
\]
where $\bar y_k$ and $\bar x_k$ are the sample means among the observations
with $z = k$.
\end{definition}

\begin{proposition}[The Wald estimator is the IV estimator]\label{prop:b3:instrumental-variables:wald}
With a binary instrument, $\hat\beta^{\mathrm{W}} = \hat\beta_1^{\mathrm{IV}}$.
\end{proposition}

\begin{proof}
Let $n_1$ observations have $z = 1$ and $n_0$ have $z = 0$, and let $f = n_1/n$,
so $\bar z = f$. For any variable $\xi$,
\[
  \sum_i (z_i - \bar z)(\xi_i - \bar\xi) = \sum_i (z_i - f) \xi_i = (1 - f)\,n_1 \bar\xi_1 - f\, n_0 \bar\xi_0
  = n f (1 - f)(\bar\xi_1 - \bar\xi_0),
\]
using $\sum_i (z_i - f) = 0$, then $(1-f) n_1 = n f (1 - f)$ and $f n_0 = n f
(1 - f)$. Apply this with $\xi = y$ and $\xi = x$; the factor $nf(1-f)$ cancels in
the ratio.
\end{proof}

\begin{example}[The Wald estimate of the return to schooling, from four means]\label{ex:b3:instrumental-variables:ak}
In the United States, a child must usually start school in the calendar year
in which she turns six and may leave on her sixteenth or seventeenth birthday.
A child born in January therefore reaches the leaving age with a little less
schooling than a child born in December. Angrist and Krueger (1991) used the
quarter of birth as an instrument for years of education. Their published
means for 327{,}509 men born 1930--1939, from the 1980 census, are:
% ledger: b3iv:ak91.panelB
\begin{center}
\begin{tabular}{lccc}
  & Born in quarter 1 & Born in quarters 2--4 & Difference (s.e.) \\ \hline
  Log weekly wage & 5.8916 & 5.9027 & $-0.01110$ (0.00274) \\
  Years of education & 12.6881 & 12.7969 & $-0.1088$ (0.0132) \\ \hline
\end{tabular}
\end{center}
The Wald estimate is $(-0.01110)/(-0.1088) = 0.1020$: one more year of
schooling raises the weekly wage by about ten log points. The authors report
it with a standard error of 0.0239, against an OLS estimate of 0.0709 (0.0003)
in the same sample. % ledger: b3iv:ak91.panelB
The ratio uses only the two differences. Every unobserved cause of wages that
is spread evenly over birth dates cancels from both.
\end{example}

The instrumental-variables estimate is larger than OLS, not smaller, though
ability bias pushes OLS up. Two readings compete. Measurement error in
reported schooling pulls OLS down, and the instrument corrects that too. Or
the men moved by compulsory schooling laws, who leave as soon as they may,
have a higher return to a year than the average man.
\Cref{sec:b3:instrumental-variables:late} makes the second reading exact.

\section{Two-stage least squares}\label{sec:b3:instrumental-variables:tsls}

Applied work rarely has one instrument and no controls. Write the model in
matrix form,
\begin{equation}\label{eq:b3:iv:matrix}
  y = X\beta + u,
\end{equation}
where $X$ is $n \times k$ and holds the constant, the exogenous controls and
the \omterm{def:b3:instrumental-variables:endogenous}{endogenous regressors}. Let $Z$ be $n \times \ell$, holding the constant,
the same exogenous controls (each is its own instrument) and the excluded
instruments, the variables that appear in $Z$ but not in $X$. The conditions
of the definition become $\E[Z'u] = 0$ (exclusion and exogeneity of the
controls) and $\operatorname{rank} \E[Z'X] = k$ (relevance), which requires
$\ell \geq k$: at least as many excluded instruments as \omterm{def:b3:instrumental-variables:endogenous}{endogenous
regressors}. With $\ell = k$ the model is just identified; with $\ell > k$ it
is over-identified. Let $P_Z = Z(Z'Z)^{-1}Z'$, the matrix that projects on the
columns of $Z$; it is symmetric and idempotent, $P_Z' = P_Z = P_Z^2$.

\begin{definition}[First stage, reduced form, two-stage least squares]\label{def:b3:instrumental-variables:tsls}
\index{first stage}\index{reduced form}\index{two-stage least squares}
In model~\eqref{eq:b3:iv:matrix} with instruments $Z$:
\begin{enumerate}
\item the \emph{first stage} is the regression of each \omterm{def:b3:instrumental-variables:endogenous}{endogenous regressor} on
all the columns of $Z$; its fitted values are $\hat X = P_Z X$ (the exogenous
columns of $X$ are their own fitted values);
\item the \emph{reduced form} is the regression of $y$ on all the columns of
$Z$;
\item the \emph{two-stage least squares} (2SLS) estimator is
\[
  \hat\beta^{\mathrm{2SLS}} = (X' P_Z X)^{-1} X' P_Z y .
\]
\end{enumerate}
\end{definition}

\begin{proposition}[What two-stage least squares does]\label{prop:b3:instrumental-variables:tsls}
\begin{enumerate}
\item $\hat\beta^{\mathrm{2SLS}}$ is the OLS coefficient of $y$ on the
first-stage fitted values $\hat X$.
\item If $\ell = k$ and $Z'X$ is invertible, $\hat\beta^{\mathrm{2SLS}} =
(Z'X)^{-1} Z'y$; with one regressor and a constant this is the estimator of
\cref{prop:b3:instrumental-variables:consistency}.
\item If $\E[Z'u] = 0$ and the rank condition holds, $\hat\beta^{\mathrm{2SLS}}$
is consistent.
\end{enumerate}
\end{proposition}

\begin{proof}
(1) OLS of $y$ on $\hat X$ is $(\hat X'\hat X)^{-1}\hat X' y$. Since $P_Z$ is
symmetric and idempotent, $\hat X'\hat X = X'P_Z'P_Z X = X'P_Z X$ and $\hat X' y =
X'P_Z y$.
(2) With $\ell = k$, $Z'X$ is square and invertible, so
$X'P_Z X = X'Z(Z'Z)^{-1}Z'X$ has the inverse $(Z'X)^{-1}(Z'Z)(X'Z)^{-1}$.
Multiply by $X'P_Z y = X'Z(Z'Z)^{-1}Z'y$: the factors $(X'Z)^{-1}X'Z$ and
$(Z'Z)(Z'Z)^{-1}$ cancel and leave $(Z'X)^{-1}Z'y$. With $Z = [\mathbf 1, z]$
and $X = [\mathbf 1, x]$, solving the two equations $Z'(y - Xb) = 0$ gives the
ratio of covariances.
(3) Substitute $y = X\beta + u$:
\[
  \hat\beta^{\mathrm{2SLS}} - \beta = \Bigl[\tfrac{X'Z}{n}\bigl(\tfrac{Z'Z}{n}\bigr)^{-1}\tfrac{Z'X}{n}\Bigr]^{-1}
  \tfrac{X'Z}{n}\bigl(\tfrac{Z'Z}{n}\bigr)^{-1}\tfrac{Z'u}{n}.
\]
Each sample moment converges to its expectation; $Z'u/n \to \E[z_i u_i] = 0$,
and the matrix in brackets converges to an invertible limit by the rank
condition. The product converges to zero.
\end{proof}

The name comes from part (1): regress each \omterm{def:b3:instrumental-variables:endogenous}{endogenous regressor} on the
instruments, then regress $y$ on the fitted values. Done literally, the second
step gives the right coefficients and the wrong standard errors, because its
residuals are $y - \hat X\hat\beta$, while the error of the model is $y -
X\beta$, with the actual regressor. Under homoskedasticity the variance is
estimated by
\[
  \widehat{\Var}\bigl(\hat\beta^{\mathrm{2SLS}}\bigr) = \hat\sigma^2 (\hat X'\hat X)^{-1},
  \qquad \hat\sigma^2 = \frac{1}{n - k}\sum_i \bigl(y_i - x_i'\hat\beta^{\mathrm{2SLS}}\bigr)^2,
\]
and the heteroskedasticity-robust version replaces $\hat\sigma^2 \hat X'\hat X$
by $\sum_i \hat u_i^2 \hat x_i \hat x_i'$ inside the sandwich of the Year 2
volume. The estimator is asymptotically normal with this variance.
\admitted
A proof follows the Year 2 volume's proof for OLS, with $\hat X$ in the place
of $X$; graduate texts on econometrics give it with the conditions on the
moments of $Z$.

\Cref{lst:b3:instrumental-variables:tsls} computes everything by hand on the
simulated town, with the college nearby as the instrument.

\omcode{code/bachelor-3/01-instrumental-variables/listings/tsls_by_hand.py}{1}{32}{The Wald ratio, \omterm{def:b3:instrumental-variables:tsls}{two-stage least squares} in matrices, the two standard errors and the first-stage $F$.}\label{lst:b3:instrumental-variables:tsls}

\omoutput{code/bachelor-3/01-instrumental-variables/outputs/tsls_by_hand.txt}{The instrument removes the ability bias; the naive second stage misstates the standard error.}\label{out:b3:instrumental-variables:tsls}

The four group means give a \omterm{def:b3:instrumental-variables:tsls}{first stage} of 0.7811 years and a \omterm{def:b3:instrumental-variables:tsls}{reduced form} of
0.0610 log points. Their ratio, 0.0781, is also the 2SLS slope, to every digit
printed, as \cref{prop:b3:instrumental-variables:tsls} says it must be; it lies
within one standard error (0.0081) of the true 0.08, where OLS gave 0.1259.
% output: tsls_by_hand, ability_bias
The naive standard error, 0.0091, is wrong by about an eighth. % output: tsls_by_hand
In real data the gap can go either way and be much larger; no applied paper
reports it, and no reader should compute 2SLS in two literal steps.

\subsection*{More instruments than regressors}

When $\ell > k$, each subset of $k$ instruments gives its own just-identified
estimate, and 2SLS combines them, weighting each instrument by how much of the
regressor it predicts. If all the instruments are valid, the estimates differ
only by sampling error. The over-identification test (Sargan's test) asks
whether they do: it regresses the 2SLS residuals on $Z$ and compares $n R^2$
with a $\chi^2_{\ell - k}$ distribution. The test can reject a set of
instruments that disagree. It can never accept the \omterm{def:b3:instrumental-variables:instrument}{exclusion restriction}: if
all the instruments share the same violation, they agree on the same wrong
number. It also assumes that every instrument estimates the same effect, which
\cref{sec:b3:instrumental-variables:late} will drop.

\subsection*{Simultaneity in supply and demand}

The oldest use of instruments is the market of the Year 1 volume. A demand
curve $q = \alpha_0 - \alpha_1 p + u^d$ and a supply curve $q = \gamma_0 +
\gamma_1 p + \delta h + u^s$, where $h$ is the weather in the growing region,
are solved together for the price and the quantity. Both depend on both
errors, so price is an \omterm{def:b3:instrumental-variables:endogenous}{endogenous regressor} in either equation. Weather moves
the supply curve and, being weather, nothing else. It is excluded from the
demand equation, and it is relevant to price. The equilibria it generates
slide along a fixed demand curve, and the Wald ratio of the change in quantity
to the change in price, between good and bad seasons, is the slope of demand,
$-\alpha_1$. A shifter of supply identifies demand; a shifter of demand would
identify supply.

\begin{omfigure}
\begin{tikzpicture}[econ, scale=0.9]
  \draw[econ axes] (0,5.2) node[above] {$p$} |- (7,0) node[right] {$q$};
  \draw[demand] (0.6,4.8) -- (6.4,0.6) node[right] {$D$};
  \draw[supply] (0.6,0.6) -- (5.4,4.8) node[above] {$S$ (bad season)};
  \draw[supply, shifted] (1.9,0.3) -- (6.6,4.4) node[right] {$S'$ (good season)};
  \coordinate (A) at (3.226,2.898);
  \coordinate (B) at (4.129,2.244);
  \node[eqpoint] at (A) {};
  \node[eqpoint] at (B) {};
  \draw[guide] (A) -- (A |- 0,0) node[below] {$q_0$};
  \draw[guide] (B) -- (B |- 0,0) node[below] {$q_1$};
  \draw[guide] (A) -- (A -| 0,0) node[left] {$p_0$};
  \draw[guide] (B) -- (B -| 0,0) node[left] {$p_1$};
  \draw[phase arrow] (3.45,2.62) -- (3.98,2.25);
\end{tikzpicture}
\par\smallskip{\small Weather shifts supply from $S$ to $S'$ and leaves demand
alone, so the two equilibria lie on the demand curve. The ratio
$(q_1 - q_0)/(p_1 - p_0)$ is its slope; a regression of $q$ on $p$ over
seasons in which both curves moved would be neither curve.}
\end{omfigure}

\section{Weak instruments}\label{sec:b3:instrumental-variables:weak}

\Cref{prop:b3:instrumental-variables:consistency} is a statement about the
limit. In a finite sample, the denominator of the IV estimator, the sample
covariance of $z$ and $x$, is itself noisy. When its true value is small, the
noise dominates.

\begin{definition}[Weak instrument]\label{def:b3:instrumental-variables:weak}
\index{weak instrument}
An instrument is a \emph{weak instrument} when the \omterm{def:b3:instrumental-variables:tsls}{first stage} explains so
little of the \omterm{def:b3:instrumental-variables:endogenous}{endogenous regressor} that the sampling distribution of the 2SLS
estimator is far from its normal approximation. In practice an instrument is
called weak when the $F$ statistic of the excluded instruments in the \omterm{def:b3:instrumental-variables:tsls}{first
stage} is small; \emph{How it is measured} below gives the thresholds.
\end{definition}

The mechanism is visible in the just-identified case. Write the \omterm{def:b3:instrumental-variables:tsls}{first stage} as
$x = \pi z + v$, where the error $v$ is correlated with $u$ (that correlation
is why $x$ is endogenous). The IV estimate minus the truth is
$\widehat{\Cov}(z, u) / (\pi \widehat{\Var}(z) + \widehat{\Cov}(z, v))$. When
$\pi$ is small, the term $\pi\widehat{\Var}(z)$ is no longer large next to the
noise $\widehat{\Cov}(z, v)$, and the noise in the numerator and the noise in
the denominator are correlated, because $u$ and $v$ are. The ratio then leans
toward $\Cov(u, v)/\Var(v)$, which is the OLS bias. A \omterm{def:b3:instrumental-variables:weak}{weak instrument} does not
merely add noise: it pulls the estimate back toward the very bias it was meant
to remove, and its usual standard error does not warn of it.
\admitted
The exact statement uses an asymptotic theory in which $\pi$ shrinks with $n$
(Staiger and Stock 1997), beyond this volume; the simulation below shows the
result.

\Cref{lst:b3:instrumental-variables:weak} sets the true effect to zero, the
correlation of $u$ and $v$ to 0.8, and $n = 200$, and chooses $\pi$ so that the
expected first-stage $F$ is 2, 10 or 50. It repeats the experiment 2{,}000
times for each.

\omcode{code/bachelor-3/01-instrumental-variables/listings/weak_iv.py}{1}{27}{2{,}000 samples at each strength of the \omterm{def:b3:instrumental-variables:tsls}{first stage}.}\label{lst:b3:instrumental-variables:weak}

\omoutput{code/bachelor-3/01-instrumental-variables/outputs/weak_iv.txt}{With $F$ near 2 the median estimate is a third of the way to OLS and a 5~\% test rejects the truth twice as often as it claims.}\label{out:b3:instrumental-variables:weak}

At an $F$ of 2 the median 2SLS estimate is 0.266, about a third of the way
from the truth, zero, to the OLS mean, 0.796; a test at 5~\% rejects the true
value in 10.7~\% of the samples. At an $F$ of 50 the median is 0.004 and the
rejection rate 4.5~\%. % output: weak_iv
With one instrument the median moves only when the instrument is very weak.
The spread is the larger damage: the chart below shows how wide the
distribution becomes.

\begin{omchart}{Sampling distribution of the 2SLS estimate when the true effect
is zero, for three strengths of the first stage (4{,}000 simulated samples of
200 each; density of the estimates in bins of 0.05). The OLS estimates of the
same samples centre near 0.77.}{simulated data, the chapter's seeded script}
\begin{tikzpicture}
\begin{axis}[omtime, xlabel={2SLS estimate (true value 0)}, ylabel={Density},
  xmin=-1, xmax=1.5, ymin=0, ymax=3.2, legend pos=north east,
  x tick label style={/pgf/number format/.cd, fixed, precision=1},
  height=0.5\linewidth]
  \addplot+[no marks, const plot mark mid] table[x=mid, y=f2] {figdata/bachelor-3/01-instrumental-variables/out/weak_iv_density.csv};
  \addplot+[no marks, const plot mark mid] table[x=mid, y=f10] {figdata/bachelor-3/01-instrumental-variables/out/weak_iv_density.csv};
  \addplot+[no marks, const plot mark mid] table[x=mid, y=f50] {figdata/bachelor-3/01-instrumental-variables/out/weak_iv_density.csv};
  \legend{$F \approx 2$, $F \approx 10$, $F \approx 50$}
\end{axis}
\end{tikzpicture}
\end{omchart}

At $F \approx 50$ the estimates gather tightly around zero. At $F \approx 2$
they spread over the whole window and beyond, with a lump between zero and the
OLS value; 27~\% of them fall outside the window altogether (against 6~\% at
$F \approx 10$), because the denominator can come close to zero. The test is the most practical worry. The second chart
runs the experiment at eleven strengths and records how often a test of
nominal size 5~\% rejects the true value.

\begin{omchart}{The size of a nominal 5~\% $t$-test of the true value after
2SLS, against the expected first-stage $F$ (4{,}000 simulated samples at each
of eleven strengths; logarithmic horizontal scale). The dotted line is
5~\%.}{simulated data, the chapter's seeded script}
\begin{tikzpicture}
\begin{semilogxaxis}[omrecord, xlabel={Expected first-stage $F$ (log scale)},
  ylabel={Share of tests rejecting the truth (\%)}, xmin=1, xmax=100, ymin=0, ymax=16,
  xtick={1,2,5,10,20,50,100}, log ticks with fixed point, height=0.45\linewidth]
  \addplot+[mark=*, mark size=1.4pt] table[x=F, y expr=100*\thisrow{reject}] {figdata/bachelor-3/01-instrumental-variables/out/weak_iv_size.csv};
  \addplot[guide, thick, domain=1:100, samples=2] {5};
\end{semilogxaxis}
\end{tikzpicture}
\end{omchart}

The rejection rate is 13.9~\% at $F = 1.5$ and 7.0~\% at $F = 10$, and it
reaches the nominal 5~\% only when $F$ is several tens.
% figdata: weak_iv_size
How fast it falls depends on the correlation between $u$ and $v$, which is
0.8 here; with a smaller correlation the test is less distorted. A reader
cannot know that correlation in real data, so the thresholds below are set for
the worst case.

\begin{measure}[The first-stage $F$ statistic]
The first-stage $F$ statistic tests that the coefficients of the excluded
instruments are all zero in the \omterm{def:b3:instrumental-variables:tsls}{first stage}; with one instrument it is the
square of that coefficient's $t$ statistic. Three thresholds are in use.
\begin{itemize}
\item \textbf{The rule of thumb.} Staiger and Stock (1997) proposed calling
instruments weak when $F < 10$. % ledger: b3iv:sy.rule
\item \textbf{Critical values.} Stock and Yogo (2005) tabulated the $F$ needed
for the actual size of a nominal 5~\% test to be at most a chosen level. With
one \omterm{def:b3:instrumental-variables:endogenous}{endogenous regressor} and one instrument, a size of at most 10~\% needs
$F > 16.38$; at most 15~\% needs $F > 8.96$; at most 20~\%, $F > 6.66$; at most 25~\%, $F > 5.53$. % ledger: b3iv:sy.crit
\item \textbf{The $tF$ procedure.} Lee, McCrary, Moreira and Porter (2022)
showed that the usual 1.96 is valid at 5~\% only when $F > 104.7$, and that a
researcher who keeps 10 as the threshold must use 3.43 in place of 1.96.
% ledger: b3iv:lmmp22.tf
\end{itemize}
With heteroskedastic or clustered errors the $F$ must itself be robust; the
classical $F$ can overstate strength. An $F$ above every threshold does not
make an instrument valid. It says only that the \omterm{def:b3:instrumental-variables:tsls}{first stage} is not the
problem.
\end{measure}

\begin{record}[Quarter of birth, examined]
Angrist and Krueger (1991) also ran regressions with many instruments: the
quarter of birth interacted with the year of birth, and with the state of
birth, up to 180 excluded instruments. Bound, Jaeger and Baker (1995)
re-examined them. With 28 instruments and controls for age within the year,
the first-stage $F$ was 1.613 and the IV estimate 0.060 (0.029), close to the
OLS estimate of 0.063; with the three quarter dummies alone the $F$ was 13.486
and the estimate 0.142 (0.033). % ledger: b3iv:bjb95.t1
They then replaced the true quarter of birth with a random draw, 500 times.
With instruments that by construction carry no information, the mean estimate
was 0.061 and the mean reported standard error 0.039, estimates that look
reasonable and are the OLS bias in disguise. % ledger: b3iv:bjb95.t3
The lesson is now part of every applied paper: report the \omterm{def:b3:instrumental-variables:tsls}{first stage}, and
add instruments only when each adds strength.
\end{record}

\section{Local average treatment effects}\label{sec:b3:instrumental-variables:late}

The model of \cref{sec:b3:instrumental-variables:endogenous} gives every
person the same $\beta_1$. Most treatments do not work that way: a year of
college pays more for some than for others, and people who expect a high
return are more likely to go. With unequal effects, what does an instrument
estimate? The answer uses the potential outcomes of the Year 2 volume, for a
binary treatment $D_i$ (going to college) and a binary instrument $Z_i$ (an
offer). Each person has two potential outcomes, $Y_i(1)$ and $Y_i(0)$, and two
potential treatments: $D_i(1)$, whether she would go if offered, and $D_i(0)$,
whether she would go if not. The observed treatment is $D_i = D_i(Z_i)$ and the
observed outcome $Y_i = Y_i(D_i)$.

\begin{definition}[Compliance types and monotonicity]\label{def:b3:instrumental-variables:compliance}
\index{complier}\index{always-taker}\index{never-taker}\index{defier}\index{monotonicity}
With a binary instrument and a binary treatment, each person has one of four
types, fixed by the pair $(D_i(0), D_i(1))$:
\begin{itemize}
\item a \emph{complier} has $D_i(0) = 0$ and $D_i(1) = 1$: treated if and only
if offered;
\item an \emph{always-taker} has $D_i(0) = D_i(1) = 1$;
\item a \emph{never-taker} has $D_i(0) = D_i(1) = 0$;
\item a \emph{defier} has $D_i(0) = 1$ and $D_i(1) = 0$: treated if and only
if not offered.
\end{itemize}
The instrument satisfies \emph{monotonicity} when there are no defiers:
$D_i(1) \geq D_i(0)$ for every $i$.
\end{definition}

\begin{omfigure}
\begin{tikzpicture}[econ, every node/.style={font=\small},
  cell/.style={draw, thick, minimum width=34mm, minimum height=12mm, align=center}]
  \node[cell, fill=omDef!10] at (0,0) {never-taker};
  \node[cell, fill=omThm!14] at (3.4,0) {complier};
  \node[cell, fill=omDebate!12] at (0,-1.2) {defier\\(ruled out)};
  \node[cell, fill=omMeth!12] at (3.4,-1.2) {always-taker};
  \node at (0,0.95) {$D_i(1) = 0$};
  \node at (3.4,0.95) {$D_i(1) = 1$};
  \node[anchor=east] at (-1.8,0) {$D_i(0) = 0$};
  \node[anchor=east] at (-1.8,-1.2) {$D_i(0) = 1$};
\end{tikzpicture}
\par\smallskip{\small The four compliance types. Columns: treated if offered?
Rows: treated if not offered? \omterm{def:b3:instrumental-variables:compliance}{Monotonicity} empties the lower-left cell.}
\end{omfigure}

The observed data never reveal a person's type: someone treated without an
offer is an \omterm{def:b3:instrumental-variables:compliance}{always-taker}, someone untreated with an offer is a \omterm{def:b3:instrumental-variables:compliance}{never-taker}, but
someone treated with an offer may be a \omterm{def:b3:instrumental-variables:compliance}{complier} or an \omterm{def:b3:instrumental-variables:compliance}{always-taker}. The
shares of the types are identified all the same.

\begin{theorem}[The local average treatment effect theorem]\label{thm:b3:instrumental-variables:late}
Suppose that
\begin{enumerate}
\item (independence) $Z_i$ is independent of $(Y_i(1), Y_i(0), D_i(1), D_i(0))$;
\item (exclusion) the instrument affects the outcome only through the
treatment, which the notation $Y_i(d)$, with no $z$ in it, already assumes;
\item (relevance) $\E[D_i \mid Z_i = 1] \neq \E[D_i \mid Z_i = 0]$;
\item (\omterm{def:b3:instrumental-variables:compliance}{monotonicity}) $D_i(1) \geq D_i(0)$ for every $i$.
\end{enumerate}
Then
\[
  \frac{\E[Y_i \mid Z_i = 1] - \E[Y_i \mid Z_i = 0]}{\E[D_i \mid Z_i = 1] - \E[D_i \mid Z_i = 0]}
  = \E\bigl[Y_i(1) - Y_i(0) \,\big|\, D_i(1) > D_i(0)\bigr],
\]
and the denominator is the share of \omterm{def:b3:instrumental-variables:compliance}{compliers}.
\end{theorem}

\begin{proof}
By exclusion, $Y_i = Y_i(0) + D_i\,(Y_i(1) - Y_i(0))$ with $D_i = D_i(Z_i)$. By
independence,
\[
  \E[Y_i \mid Z_i = 1] = \E\bigl[Y_i(0) + D_i(1)(Y_i(1) - Y_i(0))\bigr],
  \qquad
  \E[Y_i \mid Z_i = 0] = \E\bigl[Y_i(0) + D_i(0)(Y_i(1) - Y_i(0))\bigr],
\]
so the numerator is $\E\bigl[(D_i(1) - D_i(0))(Y_i(1) - Y_i(0))\bigr]$. The
difference $D_i(1) - D_i(0)$ is 1 for \omterm{def:b3:instrumental-variables:compliance}{compliers}, $-1$ for \omterm{def:b3:instrumental-variables:compliance}{defiers} and 0 for
everyone else. \omterm{def:b3:instrumental-variables:compliance}{Monotonicity} rules out \omterm{def:b3:instrumental-variables:compliance}{defiers}, so the numerator equals
$\P(\text{complier})\,\E[Y_i(1) - Y_i(0) \mid \text{complier}]$. By the same
argument, the denominator is $\E[D_i(1) - D_i(0)] = \P(\text{complier})$, which
is not zero by relevance. Divide.
\end{proof}

\begin{definition}[Local average treatment effect]\label{def:b3:instrumental-variables:late}
\index{local average treatment effect}
The \emph{local average treatment effect} (LATE) of a binary instrument is the
average treatment effect among its \omterm{def:b3:instrumental-variables:compliance}{compliers}, $\E[Y_i(1) - Y_i(0) \mid D_i(1)
> D_i(0)]$.
\end{definition}

\begin{corollary}[The shares of the types]\label{cor:b3:instrumental-variables:shares}
Under the assumptions of \cref{thm:b3:instrumental-variables:late}, the share
of \omterm{def:b3:instrumental-variables:compliance}{always-takers} is $\E[D_i \mid Z_i = 0]$, the share of \omterm{def:b3:instrumental-variables:compliance}{never-takers} is
$1 - \E[D_i \mid Z_i = 1]$, and the share of \omterm{def:b3:instrumental-variables:compliance}{compliers} is the \omterm{def:b3:instrumental-variables:tsls}{first stage}.
\end{corollary}

\begin{proof}
With no \omterm{def:b3:instrumental-variables:compliance}{defiers}, the people treated without an offer are exactly the
\omterm{def:b3:instrumental-variables:compliance}{always-takers}, and by independence their share in the group $Z_i = 0$ is their
share in the population. Likewise the people untreated with an offer are
exactly the \omterm{def:b3:instrumental-variables:compliance}{never-takers}. The three shares sum to one.
\end{proof}

\begin{model}{The local average treatment effect}
\assumes An instrument assigned as if at random, which moves the treatment in
one direction only and touches the outcome only through it; effects may differ
from person to person.
\predicts The Wald ratio equals the average effect among \omterm{def:b3:instrumental-variables:compliance}{compliers}. Two valid
instruments that move different people may give different, equally correct,
estimates. The \omterm{def:b3:instrumental-variables:tsls}{first stage} measures the share of \omterm{def:b3:instrumental-variables:compliance}{compliers}.
\tested In simulations with known effects, such as
\cref{lst:b3:instrumental-variables:late}, and in lottery and randomised-offer
studies, where independence holds by design: there the shares of the types and
the characteristics of the \omterm{def:b3:instrumental-variables:compliance}{compliers} can be computed and compared across sites
and years.
\breaks When some people do the opposite of what the instrument pushes
(\omterm{def:b3:instrumental-variables:compliance}{defiers}), the ratio mixes effects with opposite signs and can lie outside the
range of every person's effect. When the policy question is about people who
are not \omterm{def:b3:instrumental-variables:compliance}{compliers}, the LATE does not answer it.
\end{model}

\Cref{lst:b3:instrumental-variables:late} builds a population with 20~\%
\omterm{def:b3:instrumental-variables:compliance}{always-takers}, whose gain from treatment is 0.30, 30~\% \omterm{def:b3:instrumental-variables:compliance}{compliers}, whose gain
is 0.15, and 50~\% \omterm{def:b3:instrumental-variables:compliance}{never-takers}, whose gain is 0.02. The \omterm{def:b3:instrumental-variables:compliance}{always-takers} also
start from a higher untreated outcome.

\omcode{code/bachelor-3/01-instrumental-variables/listings/late_sim.py}{1}{22}{Unequal gains and three compliance types; the offer is random.}\label{lst:b3:instrumental-variables:late}

\omoutput{code/bachelor-3/01-instrumental-variables/outputs/late_sim.txt}{The Wald ratio recovers the \omterm{def:b3:instrumental-variables:compliance}{compliers}' gain, not the average gain.}\label{out:b3:instrumental-variables:late}

The treated share is 0.499 with an offer and 0.198 without, so the \omterm{def:b3:instrumental-variables:tsls}{first stage},
0.301, is the share of \omterm{def:b3:instrumental-variables:compliance}{compliers} and the share treated without an offer is the
share of \omterm{def:b3:instrumental-variables:compliance}{always-takers}. The Wald estimate, 0.143, is within sampling error of
the \omterm{def:b3:instrumental-variables:compliance}{compliers}' mean gain, 0.150; the average gain over everyone is 0.115, and
the naive comparison of treated and untreated, 0.351, adds the \omterm{def:b3:instrumental-variables:compliance}{always-takers}'
head start to their larger gain. % output: late_sim
Nothing in the data distinguishes the 0.143 from an average effect: the
assumptions do.

The theorem changes how the Wald estimate of
\cref{ex:b3:instrumental-variables:ak} reads. Its \omterm{def:b3:instrumental-variables:compliance}{compliers} are men whom the
school-leaving age kept in school for a few extra months, men who would have
left at the first chance. Their return to a year, about ten per cent, need not
be the return for a man choosing between college and work. An instrument
answers the question it was given, about the people it moves.

\section{Instruments examined}\label{sec:b3:instrumental-variables:examined}

Every instrument in applied economics has been attacked on its \omterm{def:b3:instrumental-variables:instrument}{exclusion
restriction}, and the attacks follow a few patterns. The method below collects
them.

\begin{method}[Examining an instrument]\label{met:b3:instrumental-variables:examine}
\begin{enumerate}
\item Report the \omterm{def:b3:instrumental-variables:tsls}{first stage}, its sign, its size and its robust $F$, and
compare the $F$ with the thresholds of \emph{How it is measured}.
\item Report the \omterm{def:b3:instrumental-variables:tsls}{reduced form}. If the instrument has no effect on the
outcome, the IV estimate is zero or undefined, whatever the regressions say.
\item State the story of the \omterm{def:b3:instrumental-variables:instrument}{exclusion restriction}: why the instrument is as
good as random, and why it has no path of its own to the outcome. List the
other channels and say which data could rule each out.
\item Test what can be tested: balance of the instrument across predetermined
characteristics, and placebo outcomes that the treatment cannot affect but the
other channels would.
\item Say who the \omterm{def:b3:instrumental-variables:compliance}{compliers} are, and whether the policy question is about
them.
\end{enumerate}
\end{method}

\subsection*{Geography and trade}

The opening chart is the \omterm{def:b3:instrumental-variables:tsls}{first stage} of a design due to Frankel and Romer
(1999). \Cref{lst:b3:instrumental-variables:geography} rebuilds it for 2019 and
carries it through. The gravity regression, over the 29{,}083 country pairs
that trade at all, finds a distance elasticity of $-1.40$ and a contiguity
coefficient of 0.96. % figdata: first_stage_stats
The 2SLS regression of log output per worker on the log trade share controls
for log population in both stages, because large countries trade less with
others for reasons of size alone, and size may affect productivity directly.

\omcodes{code/bachelor-3/01-instrumental-variables/listings/geography_iv.py}{1-20,21-37}{The geographic instrument, 2019: the gravity fit (lines 1--20), then the two regressions (lines 21--37).}\label{lst:b3:instrumental-variables:geography}

\omoutput{code/bachelor-3/01-instrumental-variables/outputs/geography_iv.txt}{Robust standard errors in parentheses; 157 countries.}\label{out:b3:instrumental-variables:geography}

With the population control the first-stage $F$ is 35.4, above every
threshold but the $tF$ one. OLS gives an elasticity of 0.866 (0.149) and 2SLS
gives 0.428 (0.382). % output: geography_iv
The point estimate halves, and its 95~\% interval, from about $-0.32$ to
$1.18$, includes both zero and the OLS value. On these data, geography cannot
tell whether trade raises output per worker. That is a result, not a failure:
it says how much the instrument can and cannot carry.

\begin{debate}[Does trade raise income?]
\begin{itemize}
\item Frankel and Romer (1999) estimated, with geography as the instrument for
trade in 1985, an effect of trade on income at least as large as OLS, though
imprecise. % ledger: b3iv:fr
\emph{Grade 3 on the evidence scale of
\cref{ch:b3:credible-empirical-work}}: a correlation with controls and a
plausible mechanism, since the \omterm{def:b3:instrumental-variables:instrument}{exclusion restriction} is disputed.
\item Rodríguez and Rodrik (2000) argued that the openness measures and
methods of this literature leave its findings open to other interpretations,
and found little evidence that open trade policies are associated with
growth; the same doubt applies to geography, which can reach income through
climate, disease and institutions as well as trade. \emph{Grade 3}. % ledger: b3iv:rr
\item The refit in this chapter, with 2019 data, finds a positive point
estimate that is not distinguishable from zero. \emph{Grade 3}.
\item Designs built on changes over time answer narrower questions. Feyrer
(2021) used the closing of the Suez Canal from 1967 to 1975, which lengthened
some countries' sea routes and not others', and Donaldson (2018) the arrival of
railways in the districts of colonial India; both find that lower trade costs
raised income (\cref{ch:b3:difference-in-differences} has the tools). % ledger: b3iv:feyrer, b3iv:donaldson
\emph{Grade 2} for the episodes they study, which are not the cross-country
question.
\end{itemize}
\end{debate}

\begin{record}[Rainfall, growth and conflict]
Miguel, Satyanath and Sergenti (2004) used rainfall as an instrument for
economic growth in 41 African countries over 1981--1999, arguing that rain
moves farm output and touches civil conflict through income alone. They found
that a negative growth shock of five percentage points raised the likelihood
of conflict in the following year by about one-half. % ledger: b3iv:mss04.conflict
The design was then examined on both of its conditions. On relevance, Miguel
and Satyanath (2011) reported that the \omterm{def:b3:instrumental-variables:tsls}{first stage}, the effect of rainfall on
growth, is weaker in the years after 2000. % ledger: b3iv:ms2011
On exclusion, Sarsons (2015) studied
Indian districts downstream of dams, where rainfall moves farm income much
less, and found that riots responded to rainfall there as much as elsewhere:
rain reaches conflict through some channel other than income. % ledger: b3iv:sarsons
A rainy day
also keeps people indoors.
\end{record}

The two records share a pattern. A design is proposed with a plausible story,
the \omterm{def:b3:instrumental-variables:tsls}{first stage} is measured, and later work finds either that the \omterm{def:b3:instrumental-variables:tsls}{first stage}
is weaker than it looked or a channel the story left out. Neither invalidates
the method. They are the method doing its work, in public.

\section{Exercises}\label{sec:b3:instrumental-variables:exercises}

\begin{exercise}[$\star$]\label{exo:b3:instrumental-variables:1}
Read the opening chart and its table.
\begin{enumerate}
\item The slope of the line is 0.235. Say in one sentence what it means for
two countries whose predicted trade shares differ by 20~\%.
\item Fiji's predicted share is 2.3~\% and its actual share 111~\%; the United
States' are 5.1~\% and 26~\%. Which of the two lies further from the line, in
logarithms?
\item Does the strength of this \omterm{def:b3:instrumental-variables:tsls}{first stage} say anything about whether
geography affects output per worker only through trade? Explain.
\end{enumerate}
\end{exercise}

\begin{exercise}[$\star$]\label{exo:b3:instrumental-variables:2}
Angrist and Krueger (1991) also report means for 247{,}199 men born 1920--1929,
from the 1970 census: the log weekly wage is 5.1484 for men born in the first
quarter and 5.1574 for the others, a difference of $-0.00898$; years of
education are 11.3996 and 11.5252, a difference of $-0.1256$.
% ledger: b3iv:ak91.panelA
\begin{enumerate}
\item Compute the Wald estimate of the return to a year of education from the
two printed differences.
\item The OLS estimate in the same sample is 0.0801. Is the Wald estimate
above or below it, and how does that compare with the 1980 sample of
\cref{ex:b3:instrumental-variables:ak}?
\end{enumerate}
\end{exercise}

\begin{exercise}[$\star$]\label{exo:b3:instrumental-variables:3}
In a population, $\ln w = \beta_0 + 0.06\, s + 0.12\, a + u$, where $s$ is
schooling, $a$ is unmeasured ability, and $u$ is uncorrelated with both. The
slope of a regression of $a$ on $s$ is 0.25.
\begin{enumerate}
\item Compute the probability limit of the OLS slope of $\ln w$ on $s$.
\item By how much, in per cent, does OLS overstate the return?
\item What would the limit be if ability raised schooling but not wages?
\end{enumerate}
\end{exercise}

\begin{exercise}[$\star$]\label{exo:b3:instrumental-variables:4}
A study uses the distance from a person's childhood home to the nearest
college as an instrument for her years of schooling.
\begin{enumerate}
\item State the \omterm{def:b3:instrumental-variables:instrument}{relevance condition} and the \omterm{def:b3:instrumental-variables:instrument}{exclusion restriction} for this
instrument in words.
\item Give two channels through which distance to a college could affect
adult wages other than through schooling.
\item For one of them, say what data would show that the channel is at
work.
\end{enumerate}
\end{exercise}

\begin{exercise}[$\star\star$]\label{exo:b3:instrumental-variables:5}
In a sample of $n = 1{,}000$, the sample moments of an instrument $z$, a
regressor $x$ and an outcome $y$ are $\widehat{\Var}(z) = 0.25$,
$\widehat{\Cov}(z, x) = 0.40$ and $\widehat{\Cov}(z, y) = 0.06$.
\begin{enumerate}
\item Compute the first-stage slope, the reduced-form slope and the IV
estimate.
\item The residuals $y_i - \hat\beta_0 - \hat\beta_1 x_i$ have variance
$\hat\sigma^2 = 0.36$. Under homoskedasticity the standard error of the IV
slope is $\sqrt{\hat\sigma^2 \widehat{\Var}(z) / (n\, \widehat{\Cov}(z,
x)^2)}$. Compute it and a 95~\% confidence interval.
\item Show that the formula in (2) is the $(2,2)$ entry of $\hat\sigma^2(\hat
X'\hat X)^{-1}$ when $\hat X = [\mathbf 1, \hat x]$.
\end{enumerate}
\end{exercise}

\begin{exercise}[$\star\star$]\label{exo:b3:instrumental-variables:6}
In \cref{ex:b3:instrumental-variables:ak}, IV (0.1020) exceeds OLS (0.0709).
\begin{enumerate}
\item Suppose recorded schooling is the true schooling plus a classical error
whose variance is 10~\% of the variance of recorded schooling, and there is no
ability bias. What does OLS converge to, if the true return is $\beta$? What
does IV with a valid instrument converge to?
\item Ability bias alone would make OLS exceed IV. Combine (1) with an
ability bias of $+0.01$: for which values of $\beta$ would OLS fall below IV?
\item Give a third reason, not about bias, why a valid IV estimate can exceed
OLS.
\end{enumerate}
\end{exercise}

\begin{exercise}[$\star\star$]\label{exo:b3:instrumental-variables:7}
A charity offered free tutoring, by lottery, to half of the applicants of a
Brenland school. Of the pupils offered, 62~\% took it up; of those not offered,
14~\% found a tutor elsewhere. The share passing the final examination was
0.53 among the pupils offered and 0.47 among the others.
\begin{enumerate}
\item Compute the \omterm{def:b3:instrumental-variables:tsls}{first stage}, the \omterm{def:b3:instrumental-variables:tsls}{reduced form} and the Wald estimate, and
state what the Wald estimate measures.
\item Compute the shares of \omterm{def:b3:instrumental-variables:compliance}{always-takers}, \omterm{def:b3:instrumental-variables:compliance}{never-takers} and \omterm{def:b3:instrumental-variables:compliance}{compliers}.
\item What assumption makes the share of \omterm{def:b3:instrumental-variables:compliance}{compliers} in (2) correct? Give an
example of a pupil who would violate it.
\end{enumerate}
\end{exercise}

\begin{exercise}[$\star\star$]\label{exo:b3:instrumental-variables:8}
Run \cref{lst:b3:instrumental-variables:geography} without the population
control: drop $\log$ population from both $X$ and $Z$.
\begin{enumerate}
\item Report the first-stage $F$, the OLS and the 2SLS elasticities with their
robust standard errors.
\item The first-stage $F$ falls by about half. Compute the correlations of log
population with the log built share and with the log actual share, and
explain why controlling for population strengthens the \omterm{def:b3:instrumental-variables:tsls}{first stage}.
\item Without the control, is the \omterm{def:b3:instrumental-variables:instrument}{exclusion restriction} more or less
plausible? Explain.
\end{enumerate}
\end{exercise}

\begin{exercise}[$\star\star$]\label{exo:b3:instrumental-variables:9}
In the rainfall design of the \emph{On the record} box, rainfall in year $t$
instruments growth in year $t$, and the outcome is conflict in year $t+1$.
\begin{enumerate}
\item Propose a placebo outcome: a variable that rainfall in year $t$ should
not predict if the \omterm{def:b3:instrumental-variables:instrument}{exclusion restriction} holds.
\item Sarsons's districts downstream of dams have a weaker \omterm{def:b3:instrumental-variables:tsls}{first stage}. If the
\omterm{def:b3:instrumental-variables:instrument}{exclusion restriction} held, what should the \omterm{def:b3:instrumental-variables:tsls}{reduced form} be in those
districts, compared with the others? What did she find, and what does it
imply?
\item Suppose rainfall affects conflict directly with a coefficient $c$ in
the \omterm{def:b3:instrumental-variables:tsls}{reduced form}, and the \omterm{def:b3:instrumental-variables:tsls}{first stage} is $\pi$. Write the probability limit of
the IV estimate, and say what happens to the bias when $\pi$ falls.
\end{enumerate}
\end{exercise}

\begin{exercise}[$\star\star\star$]\label{exo:b3:instrumental-variables:10}
Adapt \cref{lst:b3:instrumental-variables:weak} to a single strength, an
expected first-stage $F$ of 5, with 4{,}000 samples.
\begin{enumerate}
\item Report the mean first-stage $F$, the median 2SLS estimate, the mean OLS
estimate and the rejection rate of the true value at 5~\%.
\item What share of the samples has a first-stage $F$ above 10? Report the
median 2SLS estimate among those samples and among the others. A literature
that publishes only studies with $F > 10$ selects on the \omterm{def:b3:instrumental-variables:tsls}{first stage}: explain
the direction in which that selection moves the published estimates.
\item Compare the rejection rate with the Stock--Yogo critical values of
\emph{How it is measured}.
\end{enumerate}
\end{exercise}

\begin{exercise}[$\star\star\star$]\label{exo:b3:instrumental-variables:11}
In model~\eqref{eq:b3:iv:matrix}, let $X = [W, x]$ and $Z = [W, z]$, where $W$
holds the constant and the exogenous controls and $x$ and $z$ are single
columns.
\begin{enumerate}
\item Show that the 2SLS coefficient on $x$ equals the ratio of the
reduced-form coefficient on $z$ to the first-stage coefficient on $z$. (Hint:
use part (1) of \cref{prop:b3:instrumental-variables:tsls}, write the
first-stage fitted value as $\hat x = W a + \hat\pi z$, and apply the
Frisch--Waugh--Lovell theorem of the Year 2 volume.)
\item Deduce that the 2SLS estimate is zero whenever the reduced-form
coefficient is zero, whatever the \omterm{def:b3:instrumental-variables:tsls}{first stage}.
\end{enumerate}
\end{exercise}

\begin{exercise}[$\star\star\star$]\label{exo:b3:instrumental-variables:12}
Drop \omterm{def:b3:instrumental-variables:compliance}{monotonicity} from \cref{thm:b3:instrumental-variables:late}, keeping the
other three assumptions.
\begin{enumerate}
\item Show that the Wald ratio equals $(\pi_c \Delta_c - \pi_d
\Delta_d)/(\pi_c - \pi_d)$, where $\pi_c$ and $\pi_d$ are the shares of
\omterm{def:b3:instrumental-variables:compliance}{compliers} and \omterm{def:b3:instrumental-variables:compliance}{defiers} and $\Delta_c$ and $\Delta_d$ their average effects.
\item With $\pi_c = 0.3$, $\Delta_c = 0.1$, $\pi_d = 0.2$ and $\Delta_d = 0.5$,
compute the Wald ratio. Every person's effect is positive: comment.
\item Give a setting where \omterm{def:b3:instrumental-variables:compliance}{defiers} are plausible.
\end{enumerate}
\end{exercise}

\section{Problem: The Brenland College Lottery}\label{sec:b3:instrumental-variables:problem}

\begin{problem}[{Weekend problem --- the Brenland college lottery}]\label{pb:b3:instrumental-variables:lottery}
Brenland (invented) opened a new college in 2010 with more applicants than
places, and allocated the places by lottery. Ten years later a survey
followed all 4{,}000 applicants: the 2{,}000 who won a place and the 2{,}000
who did not. Some losers enrolled later through a waiting list or another
route; some winners never took up their place. Earnings at age thirty are in
logarithms of dalens per year.

\begin{center}
\begin{tabular}{lccccc}
  & & & \multicolumn{3}{c}{Mean log earnings} \\
  & Applicants & Enrolled & all & enrolled & not enrolled \\ \hline
  Lottery winners & 2{,}000 & 1{,}400 & 10.095 & 10.2214 & 9.80 \\
  Lottery losers & 2{,}000 & 400 & 9.995 & 10.40 & 9.89375 \\ \hline
\end{tabular}
\end{center}
The standard deviation of log earnings is 0.60 within each lottery group.

\medskip\noindent\textbf{Part I --- the naive comparison.}
\begin{enumerate}
\item Compute the share enrolled among all applicants and the mean log
earnings of all enrolled and of all non-enrolled applicants.
\item Compute the naive difference in mean log earnings between enrolled and
non-enrolled applicants.
\item Name two reasons why this difference need not be the effect of college.
\item Would a regression of log earnings on enrolment and the applicants'
school grades remove the bias? Explain in two sentences.
\item The college's prospectus quotes the naive difference as the gain from
studying there. Convert it to a percentage gain in earnings.
\end{enumerate}

\medskip\noindent\textbf{Part II --- the lottery as an instrument.}
\begin{enumerate}[resume]
\item State the \omterm{def:b3:instrumental-variables:instrument}{relevance condition} and the \omterm{def:b3:instrumental-variables:instrument}{exclusion restriction} for the
lottery, and say why the lottery makes independence credible.
\item Compute the \omterm{def:b3:instrumental-variables:tsls}{first stage} and the \omterm{def:b3:instrumental-variables:tsls}{reduced form}.
\item Compute the Wald estimate.
\item The standard error of the \omterm{def:b3:instrumental-variables:tsls}{first stage} is
$\sqrt{0.7 \times 0.3 / 2000 + 0.2 \times 0.8 / 2000}$. Compute it and the
first-stage $F$. Is the instrument weak?
\item Compute the standard error of the \omterm{def:b3:instrumental-variables:tsls}{reduced form} from the standard
deviation of log earnings within groups.
\item When the \omterm{def:b3:instrumental-variables:tsls}{first stage} is precisely estimated, the standard error of the
Wald estimate is approximately the standard error of the \omterm{def:b3:instrumental-variables:tsls}{reduced form} divided
by the \omterm{def:b3:instrumental-variables:tsls}{first stage}. Compute it.
\item Give the 95~\% confidence interval of the Wald estimate.
\end{enumerate}

\medskip\noindent\textbf{Part III --- who the compliers are.}
\begin{enumerate}[resume]
\item Assuming \omterm{def:b3:instrumental-variables:compliance}{monotonicity}, compute the shares of \omterm{def:b3:instrumental-variables:compliance}{always-takers}, \omterm{def:b3:instrumental-variables:compliance}{never-takers}
and \omterm{def:b3:instrumental-variables:compliance}{compliers}.
\item What is the mean log earnings of \omterm{def:b3:instrumental-variables:compliance}{always-takers}? Of \omterm{def:b3:instrumental-variables:compliance}{never-takers}?
\item The enrolled winners are a mix of \omterm{def:b3:instrumental-variables:compliance}{always-takers} and \omterm{def:b3:instrumental-variables:compliance}{compliers}. Use the
shares of Question~13 to show that the \omterm{def:b3:instrumental-variables:compliance}{compliers}' mean of $Y(1)$ is
$(0.7 \times 10.2214 - 0.2 \times 10.40)/0.5$, and compute it.
\item In the same way, compute the \omterm{def:b3:instrumental-variables:compliance}{compliers}' mean of $Y(0)$ from the
non-enrolled losers and the non-enrolled winners.
\item Check that the difference of the answers to Questions~15 and~16 equals
the Wald estimate.
\item Compare the \omterm{def:b3:instrumental-variables:compliance}{compliers}' mean of $Y(0)$ with the \omterm{def:b3:instrumental-variables:compliance}{never-takers}' mean, and
the \omterm{def:b3:instrumental-variables:compliance}{compliers}' mean of $Y(1)$ with the \omterm{def:b3:instrumental-variables:compliance}{always-takers}' mean. What do the two
comparisons say about who enrols, and about the naive comparison of Part I?
\end{enumerate}

\medskip\noindent\textbf{Part IV --- a second campus.}
\begin{enumerate}[resume]
\item A smaller campus held its own lottery among 1{,}200 applicants, 600
winners and 600 losers. Of the winners 26~\% enrolled, of the losers 20~\%.
Compute its \omterm{def:b3:instrumental-variables:tsls}{first stage} and the standard error
$\sqrt{0.26 \times 0.74/600 + 0.2 \times 0.8/600}$, and its first-stage $F$.
\item The \omterm{def:b3:instrumental-variables:tsls}{reduced form} at the second campus is 0.02. Compute its Wald estimate.
\item With the same within-group standard deviation of 0.60, compute the
standard error of the \omterm{def:b3:instrumental-variables:tsls}{reduced form} and, by the approximation of Question~11,
of the Wald estimate.
\item Is the second campus's estimate statistically different from the first
campus's? Compute the standard error of the difference, treating the two as
independent.
\item By the \emph{How it is measured} thresholds, how would you describe the
second instrument? Should a study pool the two lotteries as two instruments in
one 2SLS regression?
\item Even if both lotteries were strong, why might their LATEs differ? State
the effect of college that this problem estimates, with its 95~\% interval,
and for whom.
\end{enumerate}
\end{problem}
